% Numerical provenance: research/output-review-2026-11-06/results.json and
% kaggle-roadmap-2026-10-31/{local-plan.json,EXECUTION-PROTOCOL.md,degron_v3/}.
\subsection{Nested Selection and Prediction Aggregation}
\textbf{Purpose.} Obtain an out-of-fold prediction for every development cell, with all model choices made away from the cells being scored.
\textbf{Flow.}
The evaluation unit is a mutation cell, while splitting and uncertainty respect protein groups.
The full-PSI development set W contains 235 peptides, 221 groups and 122,412 cells. Five fixed outer
folds are shared across recipes. Seeds 17, 29 and 43 vary the inner group split and stochastic fitting;
they do not provide three independent outer test partitions. Arm B permits training examples with
WT PSI at most 4; arm D permits the full training range. Inner validation and outer test remain
full-range in both arms. The earlier Kaggle arm E balances WT-PSI strata within training.

All preprocessing is fitted on permitted training examples. Hyperparameters, family choices and
neural epochs are selected by inner validation MSE. A fresh model is then fitted on permitted outer
training data before predicting the untouched outer test. The nested family selector chooses among
36 frozen-feature recipes separately for each seed/fold; it does not choose the best outer OOF score.
Figure~\ref{fig:campaign-nested} distinguishes this selection from the locked transfer evaluation.

\begin{figure}[htbp]
  \centering
  % Group-aware selection and evaluation (plan-report V-2, redrawn 18 Nov 2026 to replace
% assets/diagrams/campaign-2026-11-06/nested-evaluation.pdf, which predates the LoRA-D refit).
% Sources: EXECUTION-PROTOCOL.md (W, five outer folds, inner group split, seeds 17/29/43, arms B/D,
% seed averaging, protein-group bootstrap), research/lora-refit-transfer-2026-11-16/PROTOCOL.md
% (refit on all W, T scored once). Counts are generated macros.
\begin{tikzpicture}[fig]
  \tikzset{db/.style={cylinder, shape border rotate=90, aspect=0.22, draw, fill=oiSky!20,
                      minimum width=2.6cm, minimum height=1.05cm, align=center, font=\scriptsize},
           st/.style={block, text width=3.0cm, minimum height=0.85cm},
           loop/.style={draw=oiBlue, dashed, rounded corners=6pt, inner sep=6pt},
           looplbl/.style={font=\scriptsize\bfseries, text=oiBlue, fill=white, inner sep=1.5pt}}
  % ---------------- data
  \node[db] (w) at (0,0) {W: \nVthreeNPeptides{} peptides\\\nVthreeNGroups{} protein groups};
  % ---------------- outer fold
  \node[st, fill=oiVermillion!12] (otest) at (-1.85,-1.85) {outer test groups\\full PSI range, untouched};
  \node[st, fill=oiSky!12] (otrain) at (1.85,-1.85) {outer training groups\\arm B: WT PSI $\le4$; arm D: full};
  \draw[flow] (w.south) -- ++(0,-0.55) -| (otest.north);
  \draw[flow] (w.south) -- ++(0,-0.55) -| (otrain.north);
  % ---------------- inner selection
  \node[st] (inner) at (1.85,-3.35) {inner group split; choose recipe and epochs by inner MSE};
  \node[st] (refit) at (1.85,-4.75) {fresh fit on the permitted outer training data};
  \draw[flow] (otrain) -- (inner);
  \draw[flow] (inner) -- (refit);
  \node[st, fill=oiYellow!25] (pred) at (-1.85,-4.75) {predict the outer test cells};
  \draw[flow] (refit) -- (pred);
  \draw[flow] (otest) -- (pred);
  % loops
  \node[loop, fit=(inner)(refit)] (innerloop) {};
  \node[looplbl, anchor=north east] at ($(innerloop.south east)+(-0.1,0.02)$) {inner loop};
  \node[loop, fit=(otest)(otrain)(inner)(refit)(pred), inner sep=12pt] (outer) {};
  \node[looplbl, anchor=south west] at ($(outer.north west)+(0.15,-0.02)$) {repeated: 5 outer folds $\times$ seeds 17, 29, 43};

  % ---------------- development estimates
  \node[st, text width=3.4cm, fill=oiGreen!12] (oof) at (-1.85,-7.15) {OOF prediction for every cell; mean of the three seeds};
  \node[st, text width=3.4cm, fill=oiGreen!12] (est) at (-1.85,-8.5) {peptide-macro metrics, protein-group bootstrap, OOF calibration};
  \draw[flow] (pred) -- (oof);
  \draw[flow] (oof) -- (est);

  % ---------------- locked transfer to T
  \node[db, fill=oiOrange!20] (t) at (6.1,0) {T: \nTPeptides{} peptides\\\nTGroups{} groups, held out};
  \node[st, text width=3.3cm, fill=oiOrange!12] (w5) at (6.1,-2.3) {W5-D, fixed in advance; refit on all W; scored once};
  \node[st, text width=3.3cm, fill=oiOrange!12] (lora) at (6.1,-4.2) {LoRA-D and frozen-D refit on all W after the protocol was locked; scored once (17 Nov)};
  \draw[flow] (t) -- (w5);
  \draw[flow] (w5) -- (lora);
  \node[note, text width=3.3cm, align=left, anchor=north] at ($(lora.south)+(0,-0.1)$)
       {T never enters selection or calibration; it was examined in earlier work, so it is not independent external validation.};
  \draw[dashedflow, oiGrey] (w.east) -- node[above, font=\tiny, text=oiGrey] {no group overlap} (t.west);
\end{tikzpicture}

  \caption[Group-aware selection and evaluation]{Group-aware selection and evaluation. Every model choice is made inside the outer
  training groups; the three seed predictions are averaged per cell before ensemble metrics. The
  held-out source T is scored only by recipes fixed beforehand: the W5-D ridge, and the LoRA-D and
  frozen-D refits after the transfer protocol was locked. Original schematic drawn from the
  execution and refit protocols; it replaces a diagram of 6 November that predated the refit.}
  \label{fig:campaign-nested}
\end{figure}

\textbf{Reported in} Sections~\ref{sec:res-fullpsi} and~\ref{sec:family-results}. \textbf{Limits.} The
seeds share one outer partition, and repeated human inspection of the same OOF results can still bias
later choices.

\subsection{Frozen Features and Matched Adaptation}\label{sec:campaign-adaptation}
\textbf{Purpose.} Answer RQ2 and RQ4: whether pLM representations help, and whether adapting the encoder beats a frozen encoder when everything else is held equal.
\textbf{Flow.}
The family experiment includes ESM2 at 8M, 35M, 150M and 650M parameters, ESM1b-650M,
one ESM1v-650M checkpoint, ESM-C-300M, the ProtT5 encoder and Ankh-base. For each substitution,
the mutant-minus-WT residue embedding at the mutation site is concatenated with its sequence mean.
Frozen features are reduced to 64 principal components using training-only transforms and combined
with either the common auxiliary features alone or additional ordered W5 context. Ridge regularisation
is chosen from $\{0.1,1,10,100,1000\}$ by inner MSE. The design has
$9\times2\times2\times3\times5=540$ outer fits.

The common sequence-derived inputs contain separate 20-dimensional one-hot blocks for WT and
replacement residue and a 20-dimensional peptide-composition vector. Thus a recipe labelled
``pLM-only'' still includes these 60 inputs. Measured WT PSI is used for labels, training eligibility,
balancing and diagnostic strata; it is not supplied as a predictor feature. Read counts, group IDs
and fold IDs are likewise excluded from the effect model's feature matrix.

The second experiment holds ESM2-35M, the group folds, auxiliary inputs, head initialisation and
training schedule common between frozen and LoRA branches. Both use a 64-unit ReLU head with
dropout 0.2. LoRA~\citep{hu2022lora} adds rank-8 adapters with scale $16/8$ to attention query and
value projections while retaining frozen pretrained weights. The frozen branch trains the head only.
This mutation regressor uses a pointwise/pooled embedding MLP, not a CNN or mCNN. The code also
supports full fine-tuning, but the completed campaign schedules only frozen and LoRA branches;
no full-fine-tuning performance claim follows from that implementation option.
Both branches use AdamW, learning rate $10^{-4}$, weight decay 0.01, batch size 8, FP32 computation
and gradient clipping at 1. Inner selection permits at most 20 epochs with patience 4, followed by
a fresh outer refit for the selected epoch count. The backbone remains in evaluation mode in both
branches so that its dropout policy is matched. There are 30 outer fits per branch.

The execution environment, numerical precision and tokenizer choices are recorded with the
campaign (Appendix~\ref{sec:research-local-guide}). They are part of each recipe, so the family
comparison is not a controlled comparison of pretraining alone.

\textbf{Reported in} Sections~\ref{sec:family-results} and~\ref{sec:lora-results}. \textbf{Limits.} Family
members differ in tokenizer, precision and size, so the family comparison is not a controlled test of
pretraining; full fine-tuning was implemented but not run.

\subsection{Uncertainty, Calibration and Transfer Boundaries}
\textbf{Purpose.} Turn point predictions into uncertainty statements, event probabilities and a first locked check on a different population.
\textbf{Flow.}
MAE is first averaged within a peptide and then across peptides. Positive-event metrics use
$\Delta\mathrm{PSI}\ge1$; negative-event metrics use $\Delta\mathrm{PSI}\le-1$ and require at least
three cells from each class within a peptide. Paired differences are bootstrapped over protein
groups with 5,000 resamples. Four matched local contrasts (two arms, MAE and positive AUROC)
are additionally reported with conservative post-hoc Bonferroni 98.75\% intervals. This adjustment
does not turn the development expansion into a preregistered independent confirmation.

Three-class probabilities represent negative, neutral and positive effects. Each outer-fold
calibrator is fitted using predictions on its inner validation groups, disjoint from the outer test;
probabilities are averaged across seeds. The multiclass Brier score is
$\sum_{c=1}^{3}(p_c-y_c)^2$ per cell, without division by the number of classes. Both cell-micro and
peptide-macro averages are reported. A null predictor uses inner-validation class frequencies;
positive-event ECE uses ten equal-width probability bins. Calibration is evaluated out of fold,
separately from final serving-model construction.

For transfer, W5-D is fixed before inspecting T, refitted on all permitted W data, and applied to
the locked scanning population with no recorded protein-group overlap with W or scanning development
data. T has a different substitution/label distribution and comes from a previously explored source
in the same laboratory. It does not provide a new independent external cohort. No LoRA model is
selected or evaluated using T in this development campaign; the later locked refit of
Section~\ref{sec:refit-protocol} scores LoRA-D on T once.

\textbf{Reported in} Sections~\ref{sec:campaign-calibration-transfer} and~\ref{sec:transfer-refit}.
\textbf{Limits.} Bootstrap intervals describe resampling of groups within W, not new assays; T is a
previously examined same-laboratory source.
